% Self-contained theorem/proof sections for integration into the parent report.
% Requires amsmath, amssymb, amsthm and a \Li macro.

\section{An explicit binomial parity reduction for every double weight}
\label{sec:double-parity}

Throughout this section the convention is
\[
 F_{a,b}(z)=\Li_{a,b}(z,1)
   =\sum_{n>m>0}\frac{z^n}{n^a m^b},\qquad
 U_n(z)=\operatorname{Re}\Li_n(z),\quad
 V_n(z)=\operatorname{Im}\Li_n(z).
\]
We take $|z|=1$, $z\ne1$, and use radial boundary values. These are also
the ordinary convergent series values: summation by parts applies to
$H_{n-1}^{(b)}/n^a$, whose total variation is finite and which tends to zero.
The result below is an elementary explicit specialization of the established
multiple-polylogarithm parity theorem. It should not be presented as a new
general parity theorem: Panzer proved the general theorem in 2017, including
explicit formulas at depths two and three; Umezawa supplied a general-depth
explicit formula in 2025 \cite{Panzer,Umezawa}. Its contribution here
is a short independently checkable formula in the manuscript's convention,
including the boundary cases with an index equal to one.

For $w\ge2$ define the parity component
\[
 E_{a,b}(z)=
 \begin{cases}
 \operatorname{Im}F_{a,b}(z),&w=a+b\text{ even},\\
 \operatorname{Re}F_{a,b}(z),&w=a+b\text{ odd},
 \end{cases}
\quad
 W_n(z)=
 \begin{cases}V_n(z),&w\text{ even},\\U_n(z),&w\text{ odd}.
 \end{cases}
\]
The choice of $W_n$ depends on the fixed total weight $w$, not on $n$.

\begin{theorem}[Explicit double parity formula]\label{thm:binomial-parity}
Let $w\ge2$, $1\le p\le w-1$, and $q=w-p$. Put
\begin{align*}
 A_p(z)&=\sum_{j=1}^{p}\binom{w-j-1}{q-1}E_{w-j,j}(z),\\
 P_p(z)&=
 \begin{cases}
 U_q(z)V_p(z),&w\text{ even},\ p\text{ odd},\\
 U_p(z)V_q(z),&w\text{ even},\ p\text{ even},\\
 -V_p(z)V_q(z),&w\text{ odd},\ p\text{ odd},\\
 U_p(z)U_q(z),&w\text{ odd},\ p\text{ even}.
 \end{cases}
\end{align*}
Then
\begin{align}
 A_p(z)={}&P_p(z)
 +\sum_{\substack{3\le j\le p\\j\text{ odd}}}
       \binom{w-j-1}{q-1}\zeta(j)W_{w-j}(z)
 -\sum_{\substack{3\le j\le q\\j\text{ odd}}}
       \binom{w-j-1}{p-1}\zeta(j)W_{w-j}(z)\notag\\
 &-\frac12\left[\binom{w-1}{q}-\binom{w-1}{p}\right]W_w(z)
 -\boldsymbol{1}_{w\text{ odd}}\frac{(-1)^p}{2}\zeta(w).
 \label{eq:parity-A}
\end{align}
Each individual double is recovered by the explicit inverse transform
\begin{equation}\label{eq:parity-inverse}
 E_{w-b,b}(z)=\sum_{p=1}^{b}(-1)^{b-p}
                  \binom{w-p-1}{b-p}A_p(z),\qquad 1\le b<w.
\end{equation}
Consequently every imaginary part of even total weight and every real part
of odd total weight in the same-argument double family is an explicit
polynomial in single polylogarithms and ordinary zeta values.
\end{theorem}

\begin{proof}
We supply the analytic identity responsible for the reduction, so that no
rank computation or numerical recognition enters the proof. Define
\[
 T_{p,q}(k)=\sum_{m=1}^{\infty}\frac1{(m+k)^p m^q},\qquad k\ge1.
\]
Partial fractions at $m=0$ and $m=-k$ give
\begin{align}
 \frac1{(m+k)^p m^q}
 ={}&\sum_{j=1}^{q}(-1)^{q-j}
       \binom{w-j-1}{p-1}\frac1{k^{w-j}m^j}\notag\\
 &+(-1)^q\sum_{j=1}^{p}
       \binom{w-j-1}{q-1}\frac1{k^{w-j}(m+k)^j}.
 \label{eq:parity-pf}
\end{align}
The two $j=1$ divergences cancel. Summing first to a finite cutoff and
then letting it tend to infinity yields
\begin{align}
 T_{p,q}(k)={}&
 \sum_{j=2}^{q}(-1)^{q-j}\binom{w-j-1}{p-1}
                \frac{\zeta(j)}{k^{w-j}}
 +(-1)^q\sum_{j=2}^{p}\binom{w-j-1}{q-1}
                \frac{\zeta(j)}{k^{w-j}}\notag\\
 &+(-1)^{q+1}\sum_{j=1}^{p}\binom{w-j-1}{q-1}
                \frac{H_k^{(j)}}{k^{w-j}}.
 \label{eq:parity-T}
\end{align}
Write
\[
 \mathcal A_p(z)=\sum_{j=1}^{p}\binom{w-j-1}{q-1}F_{w-j,j}(z),
 \qquad C_p=\binom{w-1}{q},
\]
and
\[
 \mathcal U_{p,q}(z)=
 \sum_{j=2}^{q}(-1)^{q-j}\binom{w-j-1}{p-1}
                 \zeta(j)\Li_{w-j}(z)
 +(-1)^q\sum_{j=2}^{p}\binom{w-j-1}{q-1}
                 \zeta(j)\Li_{w-j}(z).
\]
Using $H_k^{(j)}=H_{k-1}^{(j)}+k^{-j}$ and the hockey-stick identity,
equation~\eqref{eq:parity-T} becomes
\begin{equation}\label{eq:parity-R}
 R_{p,q}(z):=\sum_{k\ge1}z^kT_{p,q}(k)
 =\mathcal U_{p,q}(z)+(-1)^{q+1}
        \bigl(\mathcal A_p(z)+C_p\Li_w(z)\bigr).
\end{equation}

There are two product identities. The familiar same-argument shuffle is
\begin{equation}\label{eq:parity-shuffle}
 \Li_p(z)\Li_q(z)=\mathcal A_p(z)+\mathcal A_q(z).
\end{equation}
The Fourier-convolution identity is
\begin{equation}\label{eq:parity-convolution}
 \Li_p(z)\Li_q(z^{-1})
       =\zeta(w)+R_{p,q}(z)+R_{q,p}(z^{-1}).
\end{equation}
For completeness, this remains valid when $p=1$ or $q=1$. Each
$\Li_p(e^{i\theta})$ belongs to $L^2$ because $\sum n^{-2p}<\infty$.
The Fourier coefficients of their product are therefore the absolutely
convergent convolutions, equal to $\zeta(w)$ at index zero and to
$T_{p,q}(k)$ and $T_{q,p}(k)$ at positive and negative indices. Both
sequences $T_{p,q}(k)$ decrease to zero. Dirichlet's test gives convergence
of the two Fourier series, uniform on closed subarcs away from $z=1$.
Abel--Poisson summation of the $L^1$ product converges to the continuous
product at every such point, proving~\eqref{eq:parity-convolution} there.
Thus the derivation has never multiplied divergent zeta series.

If $w$ is even, take imaginary parts in
\eqref{eq:parity-shuffle}--\eqref{eq:parity-convolution}. Since $p$ and $q$
have the same parity, these give the sum and difference of $A_p,A_q$.
Moreover,
\begin{align*}
 &\operatorname{Im}\bigl(\mathcal U_{p,q}(z)+
                         \mathcal U_{q,p}(z^{-1})\bigr)\\
 &\quad=2(-1)^p\left[
  \sum_{\substack{3\le j\le p\\j\text{ odd}}}
       \binom{w-j-1}{q-1}\zeta(j)V_{w-j}(z)
 -\sum_{\substack{3\le j\le q\\j\text{ odd}}}
       \binom{w-j-1}{p-1}\zeta(j)V_{w-j}(z)\right].
\end{align*}
Solving the two elementary equations gives~\eqref{eq:parity-A} for even
$w$. If $w$ is odd, take real parts. Now $p,q$ have opposite parity and
the same two equations again give the sum and difference. The analogous
last display has $-2(-1)^p$ in place of $2(-1)^p$, with $U$ in place of
$V$. The constant term $\zeta(w)$ contributes
$-(-1)^p\zeta(w)/2$. This proves~\eqref{eq:parity-A} in the odd case too.

Finally the coefficient matrix in the definition of $A_p$ is lower
triangular with diagonal one. Its inverse follows from
\[
 \binom{w-j-1}{p-j}\binom{w-p-1}{b-p}
 =\binom{w-j-1}{b-j}\binom{b-j}{p-j}
\]
and $\sum_{p=j}^{b}(-1)^{b-p}\binom{b-j}{p-j}=\delta_{b,j}$.
This proves~\eqref{eq:parity-inverse}.
\end{proof}

\subsection{Consequences for Gaussian and Eisenstein values}
\label{subsec:gaussian-doubles}

At $z=i$ we have $V_n(i)=\beta(n)$,
$U_1(i)=-\log2/2$, and
\[
 U_n(i)=-2^{-n}(1-2^{1-n})\zeta(n),\qquad n\ge2.
\]
The odd beta values are rational multiples of powers of $\pi$.
Thus Theorem~\ref{thm:binomial-parity} supplies an explicit all-weight
reduction to the Gaussian single-value algebra in the appropriate parity.
It proves all five previously numerical weight-six equations:
\begin{align*}
2048\operatorname{Im}\Li_{5,1}(i,1)
 &=-64\pi^3\zeta(3)-527\pi\zeta(5)+4096\beta(6),\\
1536\operatorname{Im}\Li_{4,2}(i,1)
 &=96\pi^3\zeta(3)-32\pi^2\beta(4)+1581\pi\zeta(5)-8448\beta(6),\\
1024\operatorname{Im}\Li_{3,3}(i,1)
 &=-3\pi^3\zeta(3)+64\pi^2\beta(4)-1581\pi\zeta(5)+4608\beta(6),\\
23040\operatorname{Im}\Li_{2,4}(i,1)
 &=-14\pi^4\beta(2)+135\pi^3\zeta(3)-1440\pi^2\beta(4)\\
 &\hspace{17mm}+23715\pi\zeta(5)-69120\beta(6),\\
92160\operatorname{Im}\Li_{1,5}(i,1)
 &=-150\pi^5\log2+56\pi^4\beta(2)-270\pi^3\zeta(3)\\
 &\hspace{17mm}+1920\pi^2\beta(4)-675\pi\zeta(5).
\end{align*}
An especially simple all-order subfamily is
\begin{equation}\label{eq:leading-index-family}
 \operatorname{Im}\Li_{2m-1,1}(z,1)
  =(m-1)V_{2m}(z)+U_{2m-1}(z)V_1(z)
   -\sum_{k=1}^{m-1}\zeta(2k+1)V_{2m-2k-1}(z),
 \qquad m\ge1.
\end{equation}
For example, this gives the further proved weight-eight identity
\[
 \operatorname{Im}\Li_{7,1}(i,1)
 =3\beta(8)-\frac{5\pi^5}{1536}\zeta(3)
             -\frac{\pi^3}{32}\zeta(5)
             -\frac{8255\pi}{32768}\zeta(7).
\]
At $z=\rho=e^{2\pi i/3}$ write $C_n=\operatorname{Im}\Li_n(\rho)$.
The same theorem applies unchanged; in particular
\[
 \operatorname{Im}\Li_{5,1}(\rho,1)
 =2C_6-\frac{2\pi^3}{81}\zeta(3)-\frac{121\pi}{486}\zeta(5),
\qquad
 \operatorname{Im}\Li_{3,1}(\rho,1)
 =C_4-\frac{13\pi}{54}\zeta(3).
\]
These statements assert reduction identities, not numerical or motivic
independence of the remaining single values.

\section{Reconciliation of the inverse-argument mixed values}
\label{sec:mixed-reconciliation}

The following results rederive conclusions already established in the
project's 8 October 2026 continuation. They are included because the
current consolidated manuscript still reports the corresponding values
as unexplained candidate directions. They must not be counted as newly
resolved conjectures in this continuation.

For $|z|=1$, $z\ne1$, the same function introduced in the Fourier proof is
precisely the inverse-argument mixed value:
\begin{equation}\label{eq:mixed-R-identification}
 R_{a,b}(z)=\Li_{a,b}(z,z^{-1}).
\end{equation}
Indeed, put $n=m+k$ in the defining series. To justify the boundary
rearrangement when an index is one, first insert a radial factor $r^n$ in
the original outer sum. The resulting positive coefficient at difference
$k$ is
$r^k\sum_{m\ge1}r^m/((m+k)^a m^b)$; it decreases in $k$ and is bounded
by $T_{a,b}(k)$. Uniform Dirichlet tail bounds and the finite initial
segment then permit $r\uparrow1$, giving~\eqref{eq:mixed-R-identification}.

Let $\mathcal M_{a,b}$ denote the imaginary part of the mixed value at even
weight and its real part at odd weight. With the notation of
Theorem~\ref{thm:binomial-parity}, the explicit reduction is
\begin{align}\label{eq:mixed-parity-explicit}
 \mathcal M_{a,b}(z)
 ={}&(-1)^{b+1}\left[
 P_a(z)+\frac12\binom{w}{a}W_w(z)
 -\sum_{\substack{2\le j\le a\\j\text{ even}}}
     \binom{w-j-1}{b-1}\zeta(j)W_{w-j}(z)\right.\notag\\
 &\left.\hspace{24mm}
 -\sum_{\substack{2\le j\le b\\j\text{ even}}}
     \binom{w-j-1}{a-1}\zeta(j)W_{w-j}(z)\right]
 -\boldsymbol{1}_{w\text{ odd}}\frac{\zeta(w)}2,
 \qquad w=a+b.
\end{align}
This follows directly by substituting~\eqref{eq:parity-A} into
\eqref{eq:parity-R}: the odd-zeta terms cancel and
$\binom{w-1}{a}+\binom{w-1}{b}=\binom{w}{a}$.
In particular,
\begin{align*}
 \operatorname{Re}\Li_{2m,1}(z,z^{-1})
 &=\frac{2m+1}{2}U_{2m+1}(z)-\frac12\zeta(2m+1)
    +U_1(z)U_{2m}(z)\\
 &\quad-\sum_{k=1}^{m}\zeta(2k)U_{2m+1-2k}(z),\\
 \operatorname{Im}\Li_{2m-1,1}(z,z^{-1})
 &=mV_{2m}(z)+U_1(z)V_{2m-1}(z)
    -\sum_{k=1}^{m-1}\zeta(2k)V_{2m-2k}(z).
\end{align*}
At the manuscript's four proposed mixed directions these give
\begin{align*}
 \operatorname{Re}\Li_{4,1}(i,-i)
 &=\frac{3\pi^4\log2}{512}+\frac{\pi^2\zeta(3)}{64}
                               -\frac{587\zeta(5)}{1024},\\
 \operatorname{Im}\Li_{5,1}(i,-i)
 &=3\beta(6)-\frac{5\pi^5\log2}{3072}
             -\frac{\pi^2\beta(4)}6-\frac{\pi^4G}{90},\\
 \operatorname{Re}\Li_{4,1}(\rho^2,\rho)
 &=\frac{2\pi^4\log3}{243}+\frac{2\pi^2\zeta(3)}{27}
                               -\frac{281\zeta(5)}{162},\\
 \operatorname{Im}\Li_{5,1}(\rho^2,\rho)
 &=-3C_6+\frac{\pi^5\log3}{729}
             +\frac{\pi^2C_4}{6}+\frac{\pi^4C_2}{90}.
\end{align*}
Both weight-five real reductions have integral relation height below
$10^4$: clearing denominators gives coefficients
$(1024,-6,-16,587)$ and $(486,-4,-36,843)$, respectively.
Consequently their reported non-reduction against the stated complete
single-value basket is not evidence of a large-coefficient obstruction;
the value generation, basis construction, or search setup requires an
audit. The rigorous formulas settle the mathematical status regardless
of that computational history.

\section{Classical reductions for every index with one two}
\label{sec:one-two}

This section uses the one-variable abbreviation
$\Li_{k_1,\ldots,k_d}(z)=\Li_{k_1,\ldots,k_d}(z,1,\ldots,1)$.
Let $\{1\}^{r}$ denote $r$ consecutive ones. Work first on $0<z<1$, and
then continue analytically along paths avoiding the cuts of the displayed
functions. In particular, the formulas hold at $z=i$ by continuation
through the upper half-plane, with the principal values used below.

\begin{theorem}[One two at arbitrary depth]\label{thm:one-two}
For $d\ge1$, putting $a=\log(1-z)$ gives
\begin{equation}\label{eq:height-one-general}
 \Li_{2,\{1\}^{d-1}}(z)
 =\zeta(d+1)+\frac{(-1)^d a^d\log z}{d!}
   +\sum_{k=0}^{d-1}\frac{(-1)^{k+1}a^k}{k!}
                       \Li_{d+1-k}(1-z).
\end{equation}
For all $r,s\ge0$, put $\alpha=-\log(1-z)$. Then
\begin{equation}\label{eq:one-two-any-position}
 \Li_{\{1\}^{r},2,\{1\}^{s}}(z)
 =\sum_{j=0}^{r}(-1)^j\binom{s+j+1}{j}
      \frac{\alpha^{r-j}}{(r-j)!}
                     \Li_{2,\{1\}^{s+j}}(z).
\end{equation}
Consequently every value with exactly one index equal to two and all other
indices equal to one reduces to classical polylogarithms and logarithms.
\end{theorem}

The height-one functional equation is classical; a modern treatment and
arithmetic analogue appear in \cite{NakamuraShiraishi}. The proof below
also records the shuffle reduction when the two occupies any position.

\begin{proof}
The iterated integral for $d$ ones is
$\Li_{\{1\}^{d}}(z)=\alpha^d/d!$. Hence
\[
 \Li_{2,\{1\}^{d-1}}(z)
 =\frac{(-1)^d}{d!}\int_0^z\frac{\log^d(1-t)}{t}\,dt.
\]
Set $s=1-t$ and integrate by parts $d$ times. The antiderivative
\[
 -\log^d s\log(1-s)
  +\sum_{j=1}^{d}(-1)^j\frac{d!}{(d-j)!}
                       \log^{d-j}s\Li_{j+1}(s)
\]
has derivative $\log^d(s)/(1-s)$. At $s=1$ its limiting value is
$(-1)^d d!\zeta(d+1)$. This gives~\eqref{eq:height-one-general} including
the constant term and its sign.

For the second identity, locate the unique zero letter in the iterated
integral word $1^r0\,1^{s+1}$. Its variable is $t$. Integrating the
$r$ outer and $s+1$ inner identical letters gives
\[
 \Li_{\{1\}^{r},2,\{1\}^{s}}(z)
 =\frac1{r!(s+1)!}\int_0^z
    \bigl(\alpha(z)-\alpha(t)\bigr)^r\alpha(t)^{s+1}\frac{dt}{t}.
\]
Expanding the first factor and using the preceding integral formula
proves~\eqref{eq:one-two-any-position}. Analytic continuation then gives
the asserted complex values.
\end{proof}

\subsection{Proof of all three weight-four Gaussian triple formulas}
\label{subsec:gaussian-triples}

Put $L=\log2$, $G=\beta(2)$, $r=(1+i)/2$, and
$\lambda_n=\operatorname{Im}\Li_n(r)$,
$\mu_n=\operatorname{Re}\Li_n(r)$. The elementary dilogarithm reflection,
Landen transformation, and trilogarithm reflection give
\begin{align}
 \mu_2&=\frac{5\pi^2}{96}-\frac{L^2}{8},&
 \lambda_2&=G-\frac{\pi L}{8},\notag\\
 \mu_3&=\frac{35}{64}\zeta(3)-\frac{5\pi^2L}{192}
                       +\frac{L^3}{48}.
 \label{eq:halfpoint-auxiliary}
\end{align}
For clarity, the trilogarithm identity used for the last formula is
\begin{align*}
 &\Li_3(z)+\Li_3(1-z)+\Li_3\!\left(-\frac{z}{1-z}\right)\\
 &\quad=\zeta(3)+\zeta(2)\log(1-z)
  -\frac12\log z\log^2(1-z)+\frac16\log^3(1-z).
\end{align*}
It follows by differentiation using the dilogarithm identities and fixing
the constant as $z\to0$; it is then continued to $z=r$, where
$1-r=\bar r$ and $-r/(1-r)=-i$. Taking real parts proves the displayed
$\mu_3$ evaluation. The first two formulas follow respectively from
$\Li_2(r)+\Li_2(\bar r)=\zeta(2)-\log r\log\bar r$ and
$\Li_2(i)+\Li_2(\bar r)=-\tfrac12\log^2(1-i)$.

To evaluate the height-one formulas at $i$, write
\[
 a=\log(1-i)=\frac L2-\frac{\pi i}{4},\qquad
 c=\log(-(1-i))=\frac L2+\frac{3\pi i}{4}.
\]
Since $(1-i)^{-1}=r$, the needed principal-branch inversion formulas are
\begin{align*}
 \Li_2(1-i)&=-\Li_2(r)-\frac{\pi^2}{6}-\frac{c^2}{2},\\
 \Li_3(1-i)&=\Li_3(r)-\frac{\pi^2c}{6}-\frac{c^3}{6},\\
 \Li_4(1-i)&=-\Li_4(r)-\frac{7\pi^4}{360}
                         -\frac{\pi^2c^2}{12}-\frac{c^4}{24}.
\end{align*}
Substitution into~\eqref{eq:height-one-general} with $d=2,3$, using
$\log i=\pi i/2$, yields
\begin{align*}
 \operatorname{Re}\Li_{2,1}(i)&=\frac{29}{64}\zeta(3)-\frac{\pi G}{4},\\
 \operatorname{Im}\Li_{2,1}(i)&=-\frac{GL}{2}
                   +\frac{\pi L^2}{32}-\lambda_3+\frac{3\pi^3}{128},\\
 \operatorname{Im}\Li_{2,1,1}(i)&=
 \lambda_4+\frac{L\lambda_3}{2}
 -\frac{G(\pi^2-4L^2)}{32}
 -\frac{\pi(8L^3+105\zeta(3))}{768}.
\end{align*}
The real-part evaluation of $\Li_{2,1}(i)$ is also printed in
Panzer \cite{Panzer}, after reversing his index convention.
The remaining two identities follow from exact shuffle relations:
\begin{align*}
 \Li_1(z)\Li_{2,1}(z)
    &=\Li_{1,2,1}(z)+3\Li_{2,1,1}(z),\\
 \Li_2(z)\Li_{1,1}(z)
    &=3\Li_{2,1,1}(z)+2\Li_{1,2,1}(z)+\Li_{1,1,2}(z).
\end{align*}
Using $\Li_1(i)=-a$, $\Li_{1,1}(i)=a^2/2$, and
$\Li_2(i)=-\pi^2/48+iG$, one obtains
\begin{align*}
 \operatorname{Im}\Li_{1,2,1}(i)
 &=-3\lambda_4-L\lambda_3
   +\frac{G(\pi^2-4L^2)}{32}
   -\frac{3\pi^3L}{256}
   +\frac{\pi(2L^3+67\zeta(3))}{128},\\
 \operatorname{Im}\Li_{1,1,2}(i)
 &=3\lambda_4+\frac{L\lambda_3}{2}
   +\frac{5\pi^3L}{192}-\frac{163\pi\zeta(3)}{256}.
\end{align*}
These are exactly the three numerical formulas in the current manuscript,
now proved. The arbitrary-depth theorem proves additional such identities
without a new integer-relation search. It does not address the independent
status of the weight-six triples $(3,1,2),(3,2,1),(4,1,1)$, which contain
more than one zero letter and lie outside this reduction mechanism.

